\section{总结与延伸}
本讲把现代 Transformer 的大量细节组织成四类问题：norm 和 residual stream 保证训练稳定，FFN/activation 决定表达和参数效率，position embedding 决定序列几何，hyperparameters 和 attention variants 则把模型质量、训练效率和推理成本连在一起。
\begin{enumerate}
    \item \textbf{共识默认值}：pre-norm、RMSNorm、SwiGLU、RoPE、no bias 已经是现代大模型常见起点。
    \item \textbf{经验证据}：很多选择来自公开模型和消融结果，而不是完备理论证明。
    \item \textbf{系统动机}：FLOPs 少不一定 runtime 少；数据移动、KV cache 和 kernel shape 经常决定真实收益。
    \item \textbf{实验原则}：保守默认值用于起步，真正 recipe 需要小规模可控实验验证。
\end{enumerate}
\subsection{拓展阅读}
\begin{itemize}
    \item Xiong et al. 2020 on pre-norm and warmup.
    \item Zhang and Sennrich 2019 on RMSNorm.
    \item Shazeer 2020 on GLU variants and SwiGLU.
    \item Su et al. 2021 on RoPE.
    \item Shazeer 2019 and Ainslie et al. 2023 on MQA/GQA.
\end{itemize}
\end{document}
